\begin{table}[t]
\centering
\small
\setlength{\tabcolsep}{4pt}
\begin{tabular}{lcccc}
\toprule
 & \multicolumn{2}{c}{structural (share)} & \multicolumn{2}{c}{functional ($\|v\|/$mean)} \\
\cmidrule(lr){2-3}\cmidrule(lr){4-5}
scheme & Type-1 & Type-2 & at T-1 & at T-2 \\
\midrule
base             & 0.37  & 0.08 & 0.38 & 0.92 \\
triangle-CPT     & 0.38  & 0.07 & 0.41 & 0.89 \\
windowed         & 0.16  & 0.25 & 0.43 & 0.80 \\
uniform-window   & 0.21  & 0.08 & 0.45 & 0.91 \\
startup          & 0.004 & 0.32 & 0.88 & 0.81 \\
persist          & 0.005 & 0.10 & 0.86 & 0.89 \\
\bottomrule
\end{tabular}
\caption{Structural $\times$ functional transition across training schemes (Qwen2.5-0.5B, streaming eval sliding a $W{=}256$ window through the sequence). \emph{Structural:} attention share on the concentrated first token (Type-1) vs.\ the sparse separators (Type-2). \emph{Functional:} relative value norm at each sink ($\|v\|/$mean; low $\Rightarrow$ no-op, after \citealp{fesser2026unifying}). The concentrated Type-1 sink carries a \emph{low} value norm ($\sim$$0.4$, no-op-like); the sparse Type-2 separators a near-normal one ($\sim$$0.8$--$0.9$, value-carrying / broadcast). As windowed and startup training drain Type-1, attention moves onto the higher-value Type-2 --- concentrated$\to$sparse and no-op$\to$broadcast in lockstep. (Caveats: the pos-0 sink is heavily attended, so despite its low value it still injects $\sim$$15\%$ of base-model output energy --- not a \emph{pure} no-op; and uniform-window reduces Type-1 without the Type-2 spread, sinking on each window's relative boundary instead.)}
\label{tab:transition}
\end{table}
